% ED26-BM-RAT2021-4-7, exact Dubois--Prade chapter, printed 305--316.
\chapter{Representing Belief Beyond Probability and Logic}
\label{sec:p15-rationality-4-7}

An accepted proposition is used as true in further reasoning; a numerical
confidence assignment need not authorize that use. Dubois and Prade's 2021
chapter organizes alternatives to the identification of belief with one
additive probability. Its central question is which representations retain
conjunctive acceptance, which distinguish ignorance from random variation,
and which describe imprecise or conflicting information. This treatment
keeps its possibility, ranking, evidence and betting constructions separate
\cite[pp.~305--316]{DuboisPrade2021RepresentingBelief}.

The historical contrast is between a Pascalian tradition of graded
confidence and a Baconian tradition of provability or accepted certainty.
The chapter recalls that early \emph{chance} concerned repeatable outcomes,
whereas \emph{probability} could concern the credibility of testimony; it
attributes the later betting account to Ramsey and de Finetti. These are
historical source attributions, not a proof that one modern formalism is the
only rational representation. Its reference to the \emph{Encyclop\'edie}
identifies the article ``Probabilit\'e'' credited to Benjamin de Langes de
Lubi\`eres; the chapter's notes also distinguish its broad confidence
terminology from Shafer's narrower plausibility function. The survey does
not distinguish \emph{epistemic} from \emph{doxastic} state in its own
set-function discussion. Those conventions do not erase the distinction
between knowledge and belief in the multi-agent systems of Part VIII.

\section{Confidence and accepted propositions}
\label{sec:p15-r47-confidence}

\begin{hbdefinition}{def:p15-r47-capacity}{The source's confidence interface}
Let \(S\ne\varnothing\) be a state space. Events \(A,B\subseteq S\)
represent propositions, with logically equivalent propositions identified.
A confidence measure, also called a normalized capacity in this setting,
is a map \(g:2^S\to[0,1]\) satisfying
\[
 g(\varnothing)=0,\qquad g(S)=1,\qquad
 A\subseteq B\Longrightarrow g(A)\leq g(B).
\]
It describes the agent's epistemic assessment, not a many-valued truth of
the event: every proposition is still true or false. Additivity on disjoint
events is an additional property of a probability measure, not a
consequence of normalization and monotonicity. Numerical evaluations can
also be replaced by comparative confidence relations; the source names
Ramsey, de Finetti, Koopman and Lewis as historical references but does not
develop their distinct ordinal representation theorems.
\end{hbdefinition}

The betting justification adds operational assumptions. A quoted price for
a unit-paying event ticket is taken to be both an acceptable buying and an
acceptable selling price; finite portfolios must avoid sure loss. The finite
representation in \cref{def:p17-finite-price-coherence} makes that
two-sided commitment explicit. Merely supplying a number labelled
``confidence'' does not supply the commitment. Conversely, the source's
observation that exact elicitation may be demanding is a modelling argument,
not a measured psychological effect with a reported sample or error bar
\cite[p.~306]{DuboisPrade2021RepresentingBelief}.

\begin{hbdefinition}{def:p15-r47-acceptance}{Adjunction, thresholds and necessity}
A collection of accepted events excludes \(\varnothing\), contains \(S\),
is upward closed under inclusion, and, if adjunction is imposed, is closed
under binary intersection. Its Boolean indicator \(N:2^S\to\{0,1\}\)
is then a normalized confidence measure with
\[
 N(A\cap B)=\min\{N(A),N(B)\}.
\]
For numerical confidence and \(\beta>0\), threshold acceptance means
\(g(A)\geq\beta\). The source's threshold postulate requires two events
accepted at that threshold to have their intersection accepted at it.
A \emph{necessity measure} is the normalized minitive capacity obtained
when this requirement holds at every positive threshold. This universal
threshold quantifier is stronger than closure at one fixed threshold.
\end{hbdefinition}

\begin{hbtheorem}{thm:p15-r47-threshold}{proved=annex:p15.r47.threshold}{Finite accepted states and the every-threshold equivalence}
For a normalized monotone \(g\), conjunction closure at every positive
threshold is equivalent to minitivity. On finite \(S\), a Boolean minitive
\(N\) has a nonempty epistemic state
\(E=\bigcap\{A:N(A)=1\}\), and \(N(A)=1\) exactly when \(E\subseteq A\).
For a real-valued minitive \(N\), define
\[
 \iota(s)=N(S\setminus\{s\}),\quad \pi(s)=1-\iota(s),\quad
 \Pi(A)=1-N(A^c),\qquad A^c=S\setminus A.
\]
Then
\[
 N(A)=\min_{s\notin A}\iota(s),\qquad
 \Pi(A)=\max_{s\in A}\pi(s),\qquad \max_{s\in S}\pi(s)=1,
\]
with empty minimum one and empty maximum zero. In particular
\(\min\{N(A),N(A^c)\}=0\), and the positive-necessity events are
conjunction-closed and consistent.
\end{hbtheorem}

The chapter connects this representation to Zadeh's possibility theory,
Shackle's potential surprise, and Cohen's Baconian probability. The potential
surprise of \(A\) increases with \(N(A^c)\), not with its probability of
occurring. Cohen's preference for dedicated arguments in establishing guilt
is a normative philosophical proposal reported by the chapter; no general
legal standard or inference from experimental frequencies is established
here. Infinite state spaces need additional intersection or attainment
conditions: binary minitivity alone does not ensure an infinite intersection
is accepted. The finiteness hypothesis above is not removed by the source's
historical parallels \cite[p.~307 and note~4]{DuboisPrade2021RepresentingBelief}.

The source also compares Savage's act-based justification of probability
with a finite-state decision-theoretic justification of qualitative
possibility. Its cited pessimistic and optimistic criteria extend Wald's
minimum/maximum approaches; leximin and leximax refinements are proposed
when those qualitative comparisons discriminate too little. The chapter
names Dubois--Prade--Sabbadin's 2001 construction, Fargier--Sabbadin's 2005
refinement and the 2009 survey. It does not print those separate decision
axiomatizations or identify either criterion with additive expected utility
\cite[p.~308]{DuboisPrade2021RepresentingBelief}.

\section{Probability thresholds and the ranking comparison}

For an equiprobable lottery with \(n\) possible winners, the event that any
specified ticket loses has probability \(1-1/n\). All those losing events
can exceed a fixed threshold less than one, while their joint intersection
is empty. This is the source's lottery example and motivation for separating
high confidence from accepted deductively closed belief. Kyburg's alternative
abandons general adjunction; the chapter contrasts it with restricting the
allowed probability distributions, not with declaring the lottery calculation
incorrect \cite[\S2.4]{DuboisPrade2021RepresentingBelief}.

\begin{hbtheorem}{thm:p15-r47-big-stepped}{proved=annex:p15.r47.big-stepped}{A specified probability family with contextual adjunction}
On \(S=\{s_1,\ldots,s_n\}\), let an additive probability have masses
\(p_i\geq0\), totaling one, with
\[
 p_i>\sum_{j>i}p_j\qquad(1\leq i<n).
\]
For every context \(C\) with \(P(C)>0\), the most highly ordered
positive-mass state in \(C\) has more than half of its conditional mass.
Consequently the events satisfying \(P(A\mid C)>P(A^c\mid C)\)
are precisely those containing that state, and they satisfy adjunction.
This is a sufficient finite big-stepped family; it is not asserted to be
the complete classification of every probability satisfying every possible
threshold acceptance convention.
\end{hbtheorem}

For example, masses \((4,2,1)/7\) have this property. They contrast with
the uniform lottery law. A context of probability zero has no ratio
conditional here. The source also mentions Snow's atomic-bound probabilities,
Leitgeb's stability account and extensions to other thresholds; these are
distinct cited presentations, not extra conditions smuggled into the
preceding theorem.

\begin{hbdefinition}{def:p15-r47-ranking}{The static finite ranking bridge}
For a finite state space, a normalized negative ranking has state values
\(\kappa(s)\in\mathbb N\cup\{\infty\}\), with minimum zero, and event
values \(\kappa(A)=\min_{s\in A}\kappa(s)\), with
\(\kappa(\varnothing)=\infty\). Its finite possibility and necessity
images are
\[
 \Pi_\kappa(A)=2^{-\kappa(A)},\qquad
 N_\kappa(A)=1-2^{-\kappa(A^c)},\qquad 2^{-\infty}=0.
\]
These formulas preserve the different roles: low negative rank is high
possibility; positive necessity of \(A\) expresses a positive negative
rank of its complement. They do not turn a ranking conditional update
into ordinary probabilistic conditioning.
\end{hbdefinition}

\begin{hblemma}{thm:p15-r47-ranking}{proved=annex:p15.r47.ranking}{Normalization, adjunction and the interpretation boundary}
The preceding transform is a normalized finite possibility/necessity pair,
with \(\Pi_\kappa(A\cup B)=\max\{\Pi_\kappa(A),\Pi_\kappa(B)\}\)
and \(N_\kappa(A\cap B)=\min\{N_\kappa(A),N_\kappa(B)\}\).
The same numerical necessity is a coherent lower probability by the
nested-focal construction of \cref{thm:p15-ip-possibility-bridge}.
That mathematical representation does not give the rank its betting
interpretation, independence relation, or conditional dynamics.
\end{hblemma}

Ranking theory is developed separately in
\cref{sec:p15-rationality-5-3}. The 2021 source's brief ranking comparison
is read with the consistent transform above. A transformed necessity can
have both a finite lower-envelope representation and a distinct ranking
interpretation; those are different claims. The source's infinitesimal-
probability comparison and its note about well-ordering concern further
presentations, not a derivation of infinite ranking completeness from this
finite formula \cite[pp.~307, 313]{DuboisPrade2021RepresentingBelief}.

\section{Belief statements and qualitative possibility}
\label{sec:p15-r47-modal}

\begin{hbdefinition}{def:p15-r47-mel}{The two-level acceptance language}
Let \(\mathcal L\) be a classical propositional language, with truth set
\([p]\subseteq S\). The MEL language is generated by atoms \(\Box p\)
with \(p\in\mathcal L\), and outer classical Boolean connectives.
It contains neither unprefixed objective formulas as its own atoms nor
nested belief modalities. A model is a nonempty epistemic state
\(E\subseteq S\), with \(E\models\Box p\) exactly when \(E\subseteq[p]\).
Its dual \(\Diamond p=\neg\Box\neg p\) tests that \(E\) intersects
\([p]\). In particular, not believing \(p\) and believing its negation
are distinct formulas.
\end{hbdefinition}

The source describes MEL through classical tautologies on its belief atoms,
normal distribution K, consistency D, and belief in each objective tautology;
modus ponens is its propositional inference. The nested introspection schemas
of KD45 do not belong to this two-level syntax. Removing that syntax is a
change of presentation, not a proof that multi-agent KD45's relations or
interaction clauses are irrelevant. The actual KD45 sheet and card are
\cref{sheet:hilbert-kd45n,entry:p08-kd45n}; MEL's model is an information set,
not a designated actual world with an indexed agent relation
\cite[\S3.1]{DuboisPrade2021RepresentingBelief}.

\begin{hbdefinition}{def:p15-r47-gpl}{Finite grades and GPL models}
Let \(L\subseteq[0,1]\) be finite, totally ordered, and contain zero and
one. The source's graded atoms \(\Box^\lambda p\), \(\lambda\in L\setminus\{0\}\),
mean \(N([p])\geq\lambda\); outer Boolean combinations are unrestricted.
Models are normalized possibility distributions with necessity values in
\(L\), equivalently possibility values in the dual scale \(1-L\).
If \(L\) is closed under \(\lambda\mapsto1-\lambda\), those scales coincide.
The uniform grid \(L=\{0,1/k,\ldots,1\}\) in the identified 2017 author
presentation has that property.

Let \(s(\lambda)\) be the greatest value of \(L\) below positive
\(\lambda\). Then
\[
 \neg\Box^\lambda p
 \quad\Longleftrightarrow\quad N([p])\leq s(\lambda)
 \quad\Longleftrightarrow\quad
 \Pi([\neg p])\geq1-s(\lambda).
\]
This states the possibility threshold explicitly. A dual symbol carrying a
threshold is not automatically the ordinary dual at the same index; the
lowest positive grade has predecessor zero, although zero is not a
primitive necessity modality. The 2017 grid presentation writes
\(N_\lambda(p)\) for the necessity atom and defines
\(\Pi_\lambda(p)=\neg N_{1-\lambda+1/k}(\neg p)\).
These are declared notation variants, not one silently combined grammar.
\end{hbdefinition}

The source reports soundness and completeness for generalized possibilistic
logic. Its proof target is Dubois, Prade and Schockaert's 2017 presentation,
whose grid, finite objective atom set, syntax and proof are separate from
the short 2021 summary. The author version's Proposition 1 and Appendix A
are the identified citation boundary
\cite[\S3.1--3.3 and Appendix~A, author version dated 9 August 2017]
{DuboisPradeSchockaert2017AuthorGPL}. The 2021 chapter retains the full
outer Boolean language, lower-grade weakening, levelwise K/D and the
stronger consistency connection to full possibility; it does not print a
complete source proof. No new calculus or source-complete rule sheet is
asserted by this chapter's semantic definition.

\begin{hbtheorem}{thm:p15-r47-finite-modal}{proved=annex:p15.r47.finite-modal}{A complete finite semantic test and the comparative encoding}
For finite \(S\) and finite \(L\), enumerate every map
\(\pi:S\to1-L\) with maximum one. Recursive evaluation of the preceding
grammar decides finite-premise satisfaction and consequence. If
\(n=|S|\) and \(r=|L|\), exactly \(r^n-(r-1)^n\) maps are enumerated.
Furthermore
\[
 N([p])>N([q])
 \quad\Longleftrightarrow\quad
 \bigvee_{\lambda\in L\setminus\{0\}}
       (\Box^\lambda p\land\neg\Box^\lambda q).
\]
This is a semantic decision/conversion result, not a replacement proof of
the cited axiomatic completeness theorem.
\end{hbtheorem}

The standard possibilistic fragment keeps only conjunctions of weighted
objective assertions \((p,\lambda)\). The richer outer language can express
ignorance about which belief an agent holds. The source also lists relative
necessity comparisons, conditional assertions encoded through
\(N([p\Rightarrow q])>N([p\Rightarrow\neg q])\), and answer-set encodings
using three grades \(\{1,\lambda,0\}\). Its cited nonmonotonic readings
need their own model-selection/encoding rules; ordinary consequence between
sets of models remains monotone. A single ordered possibility model does
not establish the full System-P representation theorem. The smooth versus
ranked distinction remains \cref{def:p08-preferential-model,
thm:p08-klm-representation,thm:p08-rational-representation}.
Similarly, a Kleene truth-functional value is not an event confidence or
a modal assertion of ignorance. The source's references to the three-valued
translation literature remain source-specific comparisons.

\section{What one probability omits}

The chapter criticizes treating a single precise distribution as a universal
description of knowledge. Its die example contrasts substantial information
that a die is fair with no information favoring any face. Either can lead
to a uniform numerical assignment, but the assignments do not record the
different evidence histories. This does not show that a uniform chance law
is internally incoherent. It shows a loss of represented information if
the same vector is asked to encode both situations.

Its state-space example partitions extraterrestrial life first into
\(\{li,\neg li\}\), then into
\(\{ali,vli,\neg li\}\), where the first two refined outcomes mean animal
life and vegetable-only life, and \(li=ali\cup vli\). Uniform assignments
give respectively \(P(li)=1/2\) and \(P(li)=2/3\); they cannot both be one
refinement-consistent law. The positive-interval example similarly shows
that uniformity in \(x\) and uniformity in \(\log x\) are different priors,
not equivalent descriptions of ignorance. These are formal counterexamples
to representation-invariant uniformity, not measured cognitive frequencies.

\begin{hblemma}{thm:p15-r47-ignorance}{proved=annex:p15.r47.ignorance}{The two source representation-dependence calculations}
The preceding life-partition probabilities differ by \(1/6\).
If \(X\) is uniform on \([a,b]\) with \(0<a<b\), then
\(Y=\log X\) has density \(e^y/(b-a)\) on
\([\log a,\log b]\), not a uniform density. A uniform law for \(Y\)
instead induces density \(1/(x\log(b/a))\) for \(X\).
\end{hblemma}

\begin{hbdefinition}{def:p15-r47-cox}{The conditional-confidence postulates actually reported}
The source considers \(g(A\mid B)\in[0,1]\) on Boolean events with
\(B\ne\varnothing\), and the postulates
\[
\begin{aligned}
 g(A\cap C\mid B)&=F(g(A\mid C\cap B),g(C\mid B))
                      &&(C\cap B\ne\varnothing),\\
 g(A^c\mid B)&=n(g(A\mid B)) &&(B\ne\varnothing).
\end{aligned}
\]
The stated regularity assumption is that \(F\) is twice differentiable
with continuous second derivative and \(n\) twice differentiable on the
relevant argument domains. The reported conclusion is isomorphism, through
a scale transformation, to conditional probability. These are the
chapter's presentation of Cox's postulates and claimed result, not a
theorem asserting that these clauses alone suffice on every finite space.
The \(n\) here is a complement transform, not the cardinality of a lottery
or a necessity measure.
\end{hbdefinition}

The source reports the limitations of Cox's original proof, Paris's
additional infinite/density assumptions, and Halpern's finite counterexample
and infinite-setting objection. Those are cited proof/version boundaries,
not a local proof that arbitrary monotone set functions satisfy or refute
Cox's hypotheses. In particular, prescribing a complement as a function
of one confidence forces a self-duality condition; a general necessity
measure need not determine its complement that way. The chapter's Ellsberg
discussion reports ambiguity-aversion motivation but supplies no participant
sample, task protocol or quantitative estimate. It is retained as that
literature attribution, not as a new empirical claim about a population
\cite[\S4.1]{DuboisPrade2021RepresentingBelief}.

\section{Unreliable testimony and belief functions}
\label{sec:p15-r47-evidence}

\begin{hbdefinition}{def:p15-r47-mass}{Random epistemic sets, belief and plausibility}
On finite nonempty \(S\), a basic assignment has
\(m:2^S\to[0,1]\), \(m(\varnothing)=0\), and total mass one.
Positive-mass subsets are focal sets. Its event functions are
\[
 \operatorname{Bel}(A)=\sum_{\varnothing\ne E\subseteq A}m(E),\qquad
 \operatorname{Pl}(A)=\sum_{E:E\cap A\ne\varnothing}m(E).
\]
The random outcome is an information set, not an additive selection of a
point within that set. For a proper nonempty testimony \(E\) believed
reliable with probability \(p\), simple support places \(m(E)=p\) and
\(m(S)=1-p\). The remainder expresses useless information, not testimony
of \(E^c\); treating an unreliable witness as necessarily lying is a
different source model.
\end{hbdefinition}

For the degenerate testimony \(E=S\), combine the two stated weights on
the same set to obtain mass one on \(S\). A positive mass on the empty
testimony is excluded by the normalized basic-assignment convention;
conflict is handled only through the separately stated combination rule.

For a proper \(A\), simple support gives belief \(p\) if \(E\subseteq A\)
and zero otherwise. \(\operatorname{Bel}(S)=1\). Belief evaluates the
probability that the random information entails \(A\); plausibility the
probability that it is compatible with \(A\). This is the source's joining
of the probable and provable, not a declaration that either function is
the probability that the actual state lies in \(A\). The exact finite
duality and probability-envelope proof is already
\cref{thm:p14-evidence-finite-measures}.

\begin{hbtheorem}{thm:p15-r47-focal}{proved=annex:p15.r47.focal}{Adjunction, consistency and all-order belief inequalities}
For the preceding masses, belief is minitive exactly when the positive
focal sets form an inclusion chain. The weaker property
\(\min\{\operatorname{Bel}(A),\operatorname{Bel}(A^c)\}=0\) for every event
holds exactly when every two positive focal sets intersect; it need not
give a nonempty intersection of all focal sets. For every finite family
\(A_1,\ldots,A_k\),
\[
 \operatorname{Bel}\left(\bigcup_i A_i\right)\geq
 \sum_{\varnothing\ne I\subseteq\{1,\ldots,k\}}
       (-1)^{|I|+1}\operatorname{Bel}\left(\bigcap_{i\in I}A_i\right).
\]
Singleton focal sets recover an additive probability. One focal set of
mass one recovers a Boolean necessity; nested focal sets recover graded
necessity. General belief functions need not preserve adjunction.
\end{hbtheorem}

\begin{hbdefinition}{def:p15-r47-combination}{Normalized independent-source combination}
For two basic assignments, define
\[
 K=\sum_{E\cap F\ne\varnothing}m_1(E)m_2(F).
\]
When \(K>0\), normalized Dempster combination sets empty mass to zero and,
for \(A\ne\varnothing\), uses
\[
 (m_1\otimes m_2)(A)=
    \frac{\sum_{E\cap F=A}m_1(E)m_2(F)}{K}.
\]
The product models independent random information sets; normalization
conditions on their intersection being nonempty. The product and that
conditioning are substantive assumptions, not consequences of merely having
two reports. If \(K=0\), this normalized rule is undefined.
\end{hbdefinition}

\begin{hbtheorem}{thm:p15-r47-combination}{proved=annex:p15.r47.combination}{The testimony calculation and total-conflict boundary}
The preceding combination is a normalized basic assignment when \(K>0\).
Two simple supports on the same proper nonempty \(E\), with reliabilities
\(p_1,p_2\), combine to simple support
\(p_1+p_2-p_1p_2\), since \(K=1\). Two mass-one disjoint focal sets
instead have \(K=0\). Reusing the same dependent testimony does not justify
the independent-source reinforcement calculation.
\end{hbtheorem}

The source notes that only separable belief functions are obtained by
combining ordinary independent simple supports, and that Smets proposed
a generalized decomposition. Those cited decomposition statements are not
a recipe for representing every mass assignment by an arbitrary guessed
set of witnesses. Its decision discussion distinguishes the pignistic
projection of Smets and Kennes, followed by expected utility, from Jaffray's
axiomatized Hurwicz-type extension.

\begin{hbtheorem}{thm:p15-r47-pignistic}{proved=annex:p15.r47.pignistic}{The finite pignistic projection and a consonant inverse}
The pignistic probability is
\[
 \operatorname{BetP}(s)=\sum_{E\ni s}\frac{m(E)}{|E|}.
\]
It equals the Shapley allocation of the cooperative game
\(A\mapsto\operatorname{Bel}(A)\). Conversely, for ordered probability
masses \(p_1\geq\cdots\geq p_n\), put \(p_{n+1}=0\),
\(E_i=\{s_1,\ldots,s_i\}\), and \(m(E_i)=i(p_i-p_{i+1})\).
These masses are nonnegative, sum to one, and have pignistic probability
\(p\). They therefore supply an explicit nested-focal necessity inducing
\(p\); they do not show that the original belief function is uniquely
recoverable from its projection. The source's stronger least-information
characterization retains its specified literature proof boundary.
\end{hbtheorem}

\section{Lower previsions and the interpretation of imprecision}
\label{sec:p15-r47-lower-previsions}

The source's alternative betting account drops the requirement that buying
and selling prices coincide. A gamble \(f:S\to\mathbb R\) records a reward
in each state. Desirable gambles include strictly positive gains and are
closed under addition and positive scaling. The lower prevision is the
supremum acceptable buying price, with upper prevision
\(UP(f)=-LP(-f)\). A supremum need not be an accepted endpoint or an attained
maximum. Joint portfolio coherence, rather than isolated price intervals,
is the operative condition in \cref{def:p15-ip-prices,def:p15-ip-coherence}.

The actual envelope theorem is \cref{thm:p15-ip-representation}:
\[
 LP(f)=\inf_{P\in\mathcal P_{\mathcal G}}E_P(f),\qquad
 UP(f)=\sup_{P\in\mathcal P_{\mathcal G}}E_P(f).
\]
In the finite version, the compatible set consists of probabilities whose
expectations dominate every assessed buying price; an empty compatible set
signals sure loss. Coherence requires each quoted lower price to equal its
envelope, not merely to leave that set nonempty. The represented object is
a set of acceptable precise extensions, not necessarily an unknown true
probability hidden somewhere in its event intervals. On all bounded gambles
the closed convex representing set is determined; restricted event
assessments need not determine it, and a raw nonclosed prior set is not
recovered merely by taking envelopes. These qualifications are the exact
current P15 interface, not a substitution of one numerical probability.

\begin{hblemma}{thm:p15-r47-price-boundaries}{proved=annex:p15.r47.price-boundaries}{Price intervals, robust preference and maxmin are different tests}
On two exclusive exhaustive outcomes, quoting buying price \(4/5\) for
each indicator creates a sure-loss portfolio, even if each individual
selling price is \(9/10\). Thus the two isolated inequalities
\(LP\leq UP\) do not replace joint avoidance of sure loss.
For a nonempty finite credal set, strict robust preference
\(LP(f-g)>0\) implies \(LP(f)>LP(g)\), but the converse can fail.
Maxmin comparison and robust dominance are therefore distinct.
\end{hblemma}

Belief functions are a special class of lower envelopes, not the entire
class. Their finite M\"obius inverse is
\[
 m(A)=\sum_{B\subseteq A}(-1)^{|A\setminus B|}\operatorname{Bel}(B).
\]
Substituting a general lower probability in this formula gives a unique
signed inverse, which need not be nonnegative. The counterexample in
\cref{thm:p14-evidence-measure-separation} prevents replacing every coherent
lower probability by a random-set belief function. The source's proposed
behavioral characterization of belief functions by Kerkvliet and Meester
is a separate cited representation, not a proof that ordinary lower-prevision
coherence alone singles out belief functions.

\begin{hblemma}{thm:p15-r47-mobius}{proved=annex:p15.r47.mobius}{Finite inversion and the belief-function test}
The displayed subset inversion is inverse to
\(b(A)=\sum_{E\subseteq A}m(E)\). For normalized \(b\) with
\(b(\varnothing)=0\), it represents a basic assignment exactly when every
inverse mass is nonnegative. This gives a finite mathematical test, not
an inference that every envelope satisfies it.
\end{hblemma}

\section{Different possibility interpretations and the source's conclusion}

Nested-focal necessity has an evidence interpretation and a coherent lower-
probability interpretation. Giles's cautious betting interpretation, the
inverse pignistic proposal, and a static transform of ordinal rankings assign
different operational meanings to the same class of mathematical set
functions. None authorizes importing the other's update or independence
conditions. The source's language about ranking-transformed necessity being
unrelated to lower probability is therefore read as an interpretive
distinction, not as denial of the finite envelope theorem above.
In the necessity-envelope interpretation, \(N(A)>0\) entails
\(N(A^c)=0\) and hence \(\Pi(A)=1\): a positive lower indicator price
has upper price one. That mathematical relation clarifies the source's
description of cautious buying versus maximal selling prices. It does not
force acceptance of a supremum/infimum endpoint or prescribe a transaction.

The likelihood proposal is different again. For observed data \(R\),
\(\operatorname{lik}(\theta)=P(R\mid\theta)\) is not
\(P(\theta\mid R)\). On a finite parameter space, requiring maxitivity
and agreement on singletons uniquely gives
\[
 \operatorname{lik}(A)=\max_{\theta\in A}P(R\mid\theta).
\]
This conditional uniqueness follows by repeated union, not by an assumption
that all sensible confidence functions must be maxitive. A common positive
scaling factor is irrelevant to relative likelihood; turning it into a
normalized possibility distribution needs a positive maximum to divide by.
For general parameter spaces use a supremum unless a maximum is attained.
The source cites the likelihood literature for that interpretation rather
than deriving a posterior without a prior
\cite[\S4.4]{DuboisPrade2021RepresentingBelief}.

\begin{hblemma}{thm:p15-r47-likelihood}{proved=annex:p15.r47.likelihood}{Finite maxitive likelihood extension and its zero case}
For finite \(\Theta\), a maxitive set function agreeing with likelihoods
on singletons must have the displayed maximum form. If their maximum is
\(M>0\), dividing singleton values by \(M\) yields a normalized
possibility distribution. If all likelihoods are zero, this normalization
is undefined and the observed data have zero likelihood under every supplied
parameter model.
\end{hblemma}

The chapter concludes with a proposed route from graded, possibly conflicting
information toward accepted conclusions through deliberation and
argumentation. It regards some Baconian representations as descriptions of
imprecise but conflict-free information, and newer theories as able to keep
both imprecision and conflict visible. This is its comparative interpretation,
not a universal empirical taxonomy of observations. The actual P18 account
distinguishes a reason's content, source reliability, dialogue obligations
and warranted belief in \cref{sec:p18-rationality-5-5}; those conditions
remain necessary before a mass-combination or confidence score becomes an
assessment of an argument.

\begin{longtable}{@{}p{.20\linewidth}p{.34\linewidth}p{.38\linewidth}@{}}
\caption{Representation choice under the source-specific assumptions.}
\label{tab:p15-r47-choice}\\
\toprule Representation & Appropriate input and calculation & Failure of an unqualified transfer \\
\midrule\endfirsthead
\caption[]{Representation choice under the source-specific assumptions (continued)}\\
\toprule Representation & Appropriate input and calculation & Failure of an unqualified transfer \\
\midrule\endhead
\midrule\multicolumn{3}{r}{\emph{continued on next page}}\\
\endfoot
\endlastfoot
Accepted/necessity belief & A consistent information state or a minitive
confidence scale; retain conjunction at every positive threshold. & A high
additive lottery probability does not establish accepted conjunction. \\
Ranking & An ordered epistemic firmness model with its own rank scale and
conditional dynamics; the finite exponential map is explicitly static. &
Its numerical necessity image does not by itself supply betting meaning
or a probabilistic update. \\
Belief/plausibility & Masses of random information sets; containment and
intersection answer different questions. & Unreliable testimony need not
support the negation; normalized fusion needs independent sources and
positive nonconflict mass. \\
Lower/upper prevision & Joint one-sided price assessments on a stated
gamble domain; retain all compatible linear extensions. & Isolated price
intervals do not guarantee coherence, and event envelopes need not recover
the gamble model. \\
Likelihood possibility & Likelihoods for observed data across supplied
models, with an explicitly maxitive extension. & No posterior or evidence
combination rule follows without its extra prior/model assumptions. \\
\bottomrule\end{longtable}

The complete chapter bibliography is retained in the accompanying intake as
source-reference identities; their presence does not claim acquisition of
every full referenced book or article. The notes and proof targets in this
treatment identify which claims are finite constructions proved here,
which use actual existing proofs, and which remain the survey's attributed
literature claims.
\nocite{R47BanerjeeDubois2014,R47BenferhatDuboisPrade1998,R47BenferhatDuboisPrade1999,
R47ChateauneufJaffray1989,R47Cheeseman1988,R47Choquet1953,R47CiucciDubois2013,
R47Cohen1977,R47Cohen1980,R47Cox1946,R47DeFinetti1937,R47Dempster1967,
R47Dubois2006,R47DuboisFargierPrade2004,R47DuboisFargierPradeSabbadin2009,
R47DuboisLangPrade1994,R47DuboisMoralPrade1997,R47DuboisPrade1988,R47DuboisPrade1998,
R47DuboisPradeSabbadin2001,R47DuboisPradeSchockaert2017,R47DuboisPradeSmets2008,
R47Edwards1972,Ellsberg1961Ambiguity,R47FargierSabbadin2005,R47Giles1982,R47Haenni2009,
R47Halpern1999a,R47Halpern1999b,Halpern2003,Hintikka1962,R47HorvitzHeckermanLanglotz1986,
R47HuberSchmidtPetri2009,R47Jaffray1989,R47Jaynes2003,R47KerkvlietMeester2018,
R47Koopman1940,KrausLehmannMagidor1990,R47Kyburg1961,R47Kyburg1997,R47KyburgTeng2012,
R47Leitgeb2014,Lewis1973,R47Lindley1982,R47Paris1994,R47Ramsey1931,R47Savage1954,
R47Shackle1949,R47Shafer1976,R47Shafer1978,R47Shapley1953,R47Smets1995,R47SmetsKennes1994,
R47Smith1961,R47Snow1999,R47Spohn1988,R47Spohn1990,R47Sugeno1977,R47Walley1991,
R47Williams2007,R47Zadeh1978}
\index{belief representation!beyond probability and logic}
\index{possibility measure}\index{necessity measure}\index{belief function}
\index{accepted belief}\index{pignistic probability}\index{ranking!static transform}
